\section{Extending PyRATE-TA}\label{sec:extending}

\subsection{Where new code goes}

\begin{description}
  \item[\code{models/}] Kinetic model definitions. A model owns its parameter
    vector layout, its bounds, and the mapping from parameters to \Cmat{}.
    Models are pure NumPy and must be importable without Qt.
  \item[\code{fit/}] The solvers, one module per engine, plus the shared driver
    that handles parameter packing, fixed/free masks and convergence reporting.
    \code{lda.py} is the regularised LDA solver;
    \code{groundstate.py} converts difference spectra to absolute ones;
    \code{constrained.py} provides the sign-mask amplitude solver.
  \item[\code{results/}] Result dataclasses and their serialisation, carrying
    the reproducibility payload of section~\ref{sec:repro}.
    \code{lda\_result.py} is the companion to \code{fit/lda.py}.
  \item[\code{plot/}] Only what PyMORGAN cannot express: residual matrices,
    concentration profiles $\mathbf{c}(t)$, \Kmat{} / compartment diagrams,
    lifetime-density maps, fit-monitor bars, kinetic traces with residuals, and
    spectra with residuals. Everything else is delegated upstream.
  \item[\code{io/}] Import and export of fit sessions --- not raw spectroscopy
    formats, which are PyMORGAN's.
  \item[\code{gui/}] The PyQt6 application (chapter~\ref{sec:gui}).
\end{description}

\subsection{Adding a model}

A model supplies a concentration matrix and the metadata the fit driver and the
result object need: the parameter names, their bounds, and the \code{modelType}
string that will reach \code{plot\_species\_spectra}. Build \Kmat{} explicitly
and reuse the shared propagator (section~\ref{sec:kmat}) rather than deriving a
closed-form solution per scheme --- one propagator is easier to test, and the
degeneracy guard then applies automatically.

Every new model needs a synthetic-recovery test: generate data from known
parameters, fit it, assert the parameters come back within tolerance.

\subsection{Adding an export to the public API}

The public API is lazy (PEP~562). A new export needs an entry in
\code{\_EXPORTS} (name $\rightarrow$ module) \emph{and} in \code{\_\_all\_\_},
and should be added to the \code{TYPE\_CHECKING} block as well so that IDEs and
type checkers still resolve it. Heavy third-party imports used by a single
function (\code{scipy.optimise}, \code{scipy.linalg}, \code{h5py}) belong
\emph{inside} that function.

\subsection{Adding a GUI panel}

Declare the widgets in \code{main\_window.ui} with Qt Designer
(\code{uv run pyrate-ta-edit-gui}), then bind them from a mixin in
\code{gui/tabs/}. Do not build widgets in Python: the layout must stay editable
in Designer. Long operations are wrapped with \code{busy\_guard} so a second
fit cannot start inside the first, and progress is reported through a callback
the GUI adapts, keeping the engine headless.

\subsection{Adding a plot}

New analysis-only figures go in \code{plot/} as a new module, exported through
\code{plot/\_\_init\_\_.py}. Every plotter reads aesthetics from
\code{pymorgan.get\_settings()} so the figures match PyMORGAN's. The module
must not import Qt; the GUI imports the plotter, not the other way around.

\subsection{Conventions}

\begin{description}
  \item[Logging, not \code{print}.] Library code uses
    \code{logger = get\_logger(\_\_name\_\_)}. \code{logger.info} is for
    messages the user should see --- fit summaries, recovered lifetimes,
    convergence status, method citations; \code{logger.debug(\dots,
    exc\_info=True)} goes inside an \code{except} block whose failure is
    recoverable, so that nothing fails silently. \code{print} is reserved for
    console entry points and explicit \code{*\_console} helpers. To debug a
    swallowed failure:
    \code{logging.getLogger("pyrate").setLevel(logging.DEBUG)}.
  \item[Method citations.] A routine implementing a published method --- a
    global-analysis formalism, a target-analysis convention, a regularisation
    scheme --- logs its reference through \code{logger.info} when it finishes.
    An inherited habit from PyMORGAN; keep it.
  \item[No placeholders.] Write fully functional code or raise
    \code{NotImplementedError}. Deprecated code moves to \code{old\_code/}
    (gitignored, and excluded from lint and packaging) rather than being
    deleted.
  \item[Matplotlib.] Use standard mathtext; never enable
    \code{text.usetex=True}.
  \item[Style.] Target Python 3.11, 100-character lines, double quotes:
    \code{uvx ruff@latest format src tests} and
    \code{uvx ruff@latest check src tests scripts examples}.
\end{description}
